\section{卡特兰数}


\textbf{卡特兰数}（Catalan numbers）是组合计数里出现频率最高的一列数：括号匹配、格路、多边形三角剖分、二叉树、加括号方式，数出来都是同一列 $1,1,2,5,14,42,\dots$。这一节先给出它的定义与递推，再用格路与反射原理证明通项公式 $C_n=\frac1{n+1}\binom{2n}{n}$，最后说明它与三角剖分、加括号的对应关系。

\textbf{本节用到的记号：} $C_n$ 为第 $n$ 个卡特兰数（约定从 $C_0=1$ 开始）；格路指每次向右 $(1,0)$ 或向上 $(0,1)$ 走一步的路径；$\binom nk$ 为二项式系数。

\subsection{从括号序列说起}


考虑 $n$ 对左右括号组成的序列。称它是\textbf{合法的}（balanced），若从左往右读时，任意前缀中左括号的个数都不少于右括号的个数；等价地，整个序列左右括号各 $n$ 个且前缀平衡条件成立。

\begin{itemize}
\item $n=1$：只有 \texttt{()} 合法，共 $1$ 个；
\item $n=2$：\texttt{(())}、\texttt{()()} 合法，而 \texttt{())(}、\texttt{)(()} 等不合法，共 $2$ 个；
\item $n=3$：合法序列有 \texttt{((()))}、\texttt{(()())}、\texttt{(())()}、\texttt{()(())}、\texttt{()()()}，共 $5$ 个。
\end{itemize}


记 $C_n$ 为 $n$ 对括号的合法序列数。直接计算可得前几项

\[
C_0,C_1,C_2,C_3,C_4,C_5,C_6=1,1,2,5,14,42,132.
\]

\textbf{递推关系。} 任何 $n\ge1$ 的合法序列都以左括号开头。设第一个左括号所匹配的右括号在第 $2k+2$ 个位置（$k=0,1,\dots,n-1$）。那么：

\begin{itemize}
\item 第一个左括号与它匹配的右括号之间的 $2k$ 个括号，本身是一个 $k$ 对括号的合法序列，有 $C_k$ 种；
\item 这个右括号之后的剩余部分，是一个 $n-1-k$ 对括号的合法序列，有 $C_{n-1-k}$ 种。
\end{itemize}


按 $k$ 分类相加：

\[
C_n=\sum_{k=0}^{n-1}C_k\,C_{n-1-k}\qquad(n\ge1),\qquad C_0=1.
\]

这个递推是卡特兰数最本质的定义；后面几种模型都会回到它。

\subsection{格路与反射原理}


\textbf{格路模型。} 从 $(0,0)$ 出发，每次向右 $(1,0)$ 或向上 $(0,1)$，走到 $(n,n)$，且要求路径\textbf{不越过对角线} $y=x$（即始终满足 $y\le x$）。这样的路径称为 Dyck 路径。

\textbf{与括号的对应。} 把"向右"读成左括号、"向上"读成右括号：走 $n$ 步向右、$n$ 步向上，对应 $n$ 对括号；$y\le x$ 等价于"左括号数不少于右括号数"。因此 Dyck 路径数也是 $C_n$。

\textbf{总路径数。} 从 $(0,0)$ 到 $(n,n)$ 共走 $2n$ 步，其中选 $n$ 步向右，故总路径数为 $\binom{2n}{n}$。

\textbf{坏路径计数（反射原理）。} 称越过对角线的路径为坏路径，即它碰到直线 $y=x+1$。对每条坏路径，取它\textbf{第一次}碰到 $y=x+1$ 的时刻，把从起点到该时刻的这一段沿直线 $y=x+1$ 反射。起点 $(0,0)$ 反射后变成 $(-1,1)$，而反射不改变步数，于是得到一条从 $(-1,1)$ 到 $(n,n)$ 的路径；反过来也唯一还原。故坏路径与"从 $(-1,1)$ 到 $(n,n)$ 的所有路径"一一对应。

从 $(-1,1)$ 到 $(n,n)$ 需要向右 $n+1$ 步、向上 $n-1$ 步，共 $2n$ 步，因此路径数为

\[
\binom{2n}{n+1}=\binom{2n}{n-1}.
\]

于是

\[
C_n=\binom{2n}{n}-\binom{2n}{n-1}
=\binom{2n}{n}\left(1-\frac{n}{n+1}\right)
=\frac{1}{n+1}\binom{2n}{n}.
\]

这就是卡特兰数的通项公式。

\textbf{例（$n=4$）。} 总路径数 $\binom84=70$，坏路径数 $\binom83=56$，好路径数 $70-56=14=C_4$。

\subsection{多边形三角剖分与加括号}


\textbf{三角剖分。} 一个凸 $(n+2)$ 边形，用 $n-1$ 条互不相交的对角线把它剖分成 $n$ 个三角形，称为一个\textbf{三角剖分}。三角剖分的个数也是 $C_n$。

\textbf{证明（按含固定边的三角形分类）。} 设多边形顶点依次为 $v_1,v_2,\dots,v_{n+2}$，固定边 $v_1v_2$。在任何三角剖分中，边 $v_1v_2$ 恰好属于某一个三角形，设这个三角形的第三个顶点为 $v_k$（$3\le k\le n+2$）。这条对角线把多边形分成两个较小的多边形：

\begin{itemize}
\item 顶点为 $v_2,v_3,\dots,v_k$ 的 $(k-1)$ 边形；
\item 顶点为 $v_1,v_k,v_{k+1},\dots,v_{n+2}$ 的 $(n+4-k)$ 边形。
\end{itemize}


前者的三角剖分数是 $C_{k-3}$（因为 $(k-1)$ 边形对应下标 $k-3$），后者是 $C_{n+2-k}$（$(n+4-k)$ 边形对应下标 $n+2-k$）。令 $a=k-3$，则 $a$ 从 $0$ 到 $n-1$，得到

\[
C_n=\sum_{a=0}^{n-1}C_a\,C_{n-1-a},
\]

正是卡特兰数的递推。$\square$

\textbf{加括号。} 给 $n+1$ 个符号的乘积加\textbf{完全括号}，不同的加括号方式数也是 $C_n$。例如 $n=2$（3 个符号 $a,b,c$）有 $(ab)c$ 与 $a(bc)$ 两种。

\textbf{为什么与三角剖分一致。} 固定 $(n+2)$ 边形的一条边 $e$ 作为"最终结果"，其余 $n+1$ 条边各看成一个符号。三角剖分中每个三角形看成一次乘法：它把三条边对应的对象合并成一个新对象；从叶子三角形开始逐层合并，最后在 $e$ 上得到整个乘积。不同的剖分给出不同的合并顺序，也就给出不同的加括号方式。反过来，给定 $n+1$ 个符号的完全加括号方式，从最内层开始把每次乘法画成一个三角形、逐层向外扩展，就唯一地得到 $(n+2)$ 边形的一个三角剖分。因此

\[
\text{三角剖分}\ \longleftrightarrow\ \text{加括号方式}\ \longleftrightarrow\ \text{合法括号序列},
\]

三者数量都等于 $C_n$。

\subsection{其他等价模型与提醒}


\textbf{等价模型。} $C_n$ 还等于：

\begin{itemize}
\item $n$ 个\textbf{内部结点}（$n+1$ 片叶子）的满二叉树的个数；
\item $n$ 对括号的合法序列数；
\item 从 $(0,0)$ 到 $(n,n)$ 不越过对角线的格路数；
\item $(n+2)$ 边形的三角剖分数；
\item $n+1$ 个符号的完全加括号方式数；
\item $n$ 个元素经一个栈得到的\textbf{合法出栈序列}数（补充，见 [DS] §11.1 例题）。
\end{itemize}


\textbf{下标不要错位。} $C_n$ 对应"$n$ 对括号 / $n+2$ 边形 / $n+1$ 个符号 / 不越过对角线的 $n\times n$ 格路"。做题时先确认题目里的 $n$ 是哪一个，再套公式。

\textbf{公式自检。} $n=0$ 时 $\frac1{1}\binom00=1=C_0$；$n=1$ 时 $\frac12\binom21=1$；$n=2$ 时 $\frac13\binom42=2$。用递推算出的前几项与公式逐项一致。

\textbf{本节的勘误与提醒。}
\begin{itemize}
\item 手写稿里"$C_{2n}=\binom{2n}{n}-\binom{2n}{n-1}$"中的 $C_{2n}$ 是笔误，左边应是第 $n$ 个卡特兰数 $C_n$；$\binom{2n}{n}-\binom{2n}{n-1}$ 是\textbf{计数结果}，不是卡特兰数的下标。
\item 反射原理中"坏路径"要按\textbf{第一次}碰到 $y=x+1$ 反射，且反射的是起点到该时刻的一段；反射后终点不变，起点变成 $(-1,1)$。
\item 括号序列的合法条件是"任意前缀左括号不少于右括号"，不是"左右括号总数相等"就够。
\item 三角剖分中固定的是\textbf{一条边}（而不是一个顶点）；分类用的三角形必须含这条边，分类才不重不漏。
\item 课堂转写把"多边形三角剖分"和"加括号"分成两段讲，且第二段开头从多边形讲起；手写稿则先画格路。两处说的是同一件事，整理时按"格路 → 递推 → 三角剖分/加括号"的顺序合并。
\end{itemize}
